%======================================================================
\section{The channel of the configuration lamps, and Theorem B}\label{sec:lamp-channel}
%======================================================================

This section builds the channel of Theorem~B from the certificate of Theorem~A and states the theorem in full.
Theorem~\ref{thm:weighted} needs $a$, $b$ and a name for $ab$ among the generators. Let $R$ be the certificate of Table~\ref{tab:certificate}, in
the five generators $p,q,t,c,a$, with $b$ standing for the word $ac^4$. Let $S_\Phi=\{p,q,t,c,a,b,j\}$, where the new letters represent
$b=ac^4$ and $j=ab$ in $\Gamma$, and let $R_\Phi$ consist of the $25$ words of $R$ and the four words
\[
b^{-1}ac^4,\qquad b^2,\qquad (ab)^3,\qquad j^{-1}ab
\]
in the letters of $S_\Phi$. They represent $1$ in $\Gamma$, $a^2\in R$, and $j\ne1$ because $ab$ has order three (Lemma~\ref{lem:structure}). The
channel of the configuration lamps is the gadget channel $\Phi=\Phi_\Gamma$ of $R_\Phi$.

\paragraph{Dimension and rank.} The channel by the numbers:
\begin{center}\small
\begin{tabular}{cccccccc}
\toprule
letters & words & total length & symbols & triangles & longest word & dimension & Choi rank\\
\midrule
$s=7$ & $r_0=29$ & $\Lambda=320$ & $276$ & $298$ & $m=24$ & $d=873$ & $r_*=130$\\
\bottomrule
\end{tabular}
\end{center}
The total length is $\Lambda=303+6+2+6+3$, and~\eqref{eq:gadget-dimension} gives $d=1+4\cdot7+3\cdot320-4\cdot29=873$. The Choi rank $r_*$ is
the number of distinct labels. Of the $277$ candidates, the identity and the $276$ symbols, exactly $130$ are distinct elements of
$\Gamma$, so $r_*=130$. Most coincidences are shared prefixes: $pq$ and $pqp^{-1}$, for instance, begin several words. The others are relations
such as $a^2=1$ and $j=ab=c^4$. The script \texttt{verification/choi\_rank\_check.py} computes this partition with two independent engines,
the faithful action of $\Gamma$ on configurations and the embedding in $3V$ (Appendix~\ref{app:machine-checks}). The rank rate is therefore
$\frac12\log130\approx3.51$ qubits per use.

\paragraph{Transport.} For $N\ge1$, the proof of Theorem~\ref{thm:main-group} provides $2^N$ words $T_u$, one for each configuration $u$
supported in the first $N$ points of ray~$1$, with $|T_u|\le21N\max(3,N)\le63N^2$ (Theorem~\ref{thm:pair-area}) and
$\Area_{R_\Phi}([T_uxT_u^{-1},T_wyT_w^{-1}])\le\Lambda_N+4\le C_1N^{\kappa}$ for $u\ne w$ and $x,y\in\{a,b\}$, where
\[
p_H=\log107+4,\qquad \kappa=2+p_H<12.75,
\]
and $C_1$ is absolute; the letter $b$ is converted to $ac^4$ by the word $b^{-1}ac^4$, at most four times, and areas with respect to
$R_\Phi\supseteq R$ are no larger than with respect to $R$. So the hypothesis of Corollary~\ref{cor:amplification}(a) holds with this $\kappa$ and
$K=C_1+64$.

\begin{theorem}[Theorem B]\label{thm:main-channel}
The channel $\Phi$ on $M_{873}$ lies in $\cK_{873}$, and its Choi rank is $r_*=130$. Let $\kappa=2+p_H=6+\log107$. There are
constants $c,c_0,K>0$ and $n_0$ such that the following hold; all services have error at most $1/16$.
\begin{enumerate}[label=(\alph*),itemsep=1pt]
\item \emph{Memory against purity.} For $n\ge n_0$ and $\log n\le B\le c_0n$, every causal service of $\Phi$ for $n$ rounds with purity
$P+C\le B$ has, at any exchange,
\[
Q\ \ge\ c\,\frac{n}{B\,(\log n)^{2\kappa}} .
\]
For every $n\ge1$ and $0<B\le n$ there is an exact causal service, with error $0$, purity $P+C<B$, exchange $J=(n-1)Q$ and memory
$Q\le K(n/B)\log(2n/B)$. In particular, for every $A_0\ge0$ there is $c'>0$ such that services with $P+C\le A_0\log n$ have
$Q\ge c'n/(\log n)^{2\kappa+1}$ for all large $n$, where $2\kappa+1=2p_H+5<26.49$, and there are services with $P=C=0$ and $Q\le Kn\log n$.
\item \emph{The rank rate.} For $n\ge n_0$, every causal service of $\Phi$ for $n$ rounds at the rank rate, $J\le\frac n2\log r_*$, has, whatever
its purity,
\[
Q\ \ge\ c\,\frac{\sqrt n}{(\log n)^{\kappa}} .
\]
For every $n\ge2$ there is a causal service at the rank rate with $C=0$ and $P,Q\le K\sqrt{n\log n}$.
\end{enumerate}
So memory and purity trade at a fixed rate: for $n\ge n_0$ and $\log n\le B\le c_0n$ the least memory of a service with purity at most $B$ lies
between $c\,n/(B(\log n)^{2\kappa})$ and $K(n/B)\log(2n/B)$, and the product of memory and purity is $n^{1+o(1)}$. At the rank rate the least
memory is $n^{1/2+o(1)}$, and it is attained with purity of the same order: the rank rate sits where memory and purity are equal.
\end{theorem}

% Figure: memory against purity for the lamp channel (schematic, logarithmic axes).
\begin{figure}[t]
\centering
\begin{tikzpicture}[x=1cm,y=1cm,font=\small,>=Stealth,lbl/.style={font=\footnotesize,inner sep=1pt}]
\colorlet{accent}{blue!50!black}
% below the band: no service (Theorem B(a)); the band: Q.B = n^{1+o(1)}, for log n <= B <= c_0 n
\fill[black!8] (1.0,5.65) -- (6.0,0.65) -- (6.0,0) -- (1.0,0) -- cycle;
\fill[accent!22] (1.0,6.35) -- (6.0,1.35) -- (6.0,0.65) -- (1.0,5.65) -- cycle;
\draw[accent] (1.0,6.35) -- (6.0,1.35);
\draw[accent] (1.0,5.65) -- (6.0,0.65);
\node[lbl,anchor=west] at (6.1,1.35) {$K(n/B)\log(2n/B)$: exact services};
\node[lbl,anchor=west] at (6.1,0.65) {$c\,n/(B(\log n)^{2\kappa})$};
\node[lbl,rotate=-45] at (2.55,4.45) {$Q\cdot B=n^{1+o(1)}$};
\node[lbl] at (3.3,1.1) {no service (Theorem~B(a))};
% the purity budget at the rank rate
\draw[dashed] (2.6,2.2) -- (6.9,6.5);
\node[lbl,anchor=west,align=left] at (6.95,6.5) {$P+C=3Q+6$: rank-rate services\\ lie to the left (purity budget)};
% axes
\draw[->] (0,0) -- (8.0,0);
\node[lbl,anchor=north east] at (8.0,-0.12) {purity $B$};
\draw[->] (0,0) -- (0,7.4);
\node[lbl,anchor=south west] at (0.1,7.3) {memory $Q$};
\draw (0.42,-0.12) -- (0.58,0.12) (0.52,-0.12) -- (0.68,0.12);
\foreach \x/\l in {0/{$0$},1.0/{$\log n$},3.5/{$\sqrt n$},6.0/{$c_0n$}}{
  \draw (\x,0.06) -- (\x,-0.06); \node[lbl,anchor=north] at (\x,-0.12) {\l};
}
\foreach \y/\l in {3.5/{$\sqrt n$},6.0/{$n$},6.6/{$n\log n$}}{
  \draw (0.06,\y) -- (-0.06,\y); \node[lbl,anchor=east] at (-0.12,\y) {\l};
}
% the marked points
\fill (0,6.6) circle (2pt);
\node[lbl,anchor=west] at (0.2,6.75) {zero purity: $Q\le Kn\log n$};
\fill (3.5,3.5) circle (2pt);
\draw (3.5,3.5) -- (4.55,3.55);
\node[lbl,anchor=west,align=left] at (4.6,3.55) {rank rate: $Q$ and $P+C$\\ both about $\sqrt n$};
\end{tikzpicture}
\caption{Memory against purity for the lamp channel, on logarithmic axes. For $\log n\le B\le c_0n$ the least memory of a service with purity at
most $B$ lies in the band $Q\cdot B=n^{1+o(1)}$, between $c\,n/(B(\log n)^{2\kappa})$ and $K(n/B)\log(2n/B)$ (Theorem~\ref{thm:main-channel}(a)),
and at zero purity memory $Kn\log n$ suffices. At the rank rate the purity is less than $3Q+6$ (Proposition~\ref{prop:purity-budget}), so the
band forces memory about $\sqrt n$, and streaming services attain this with purity of the same order (Theorem~\ref{thm:main-channel}(b)).}
\label{fig:tradeoff}
\end{figure}


The exponent $2\kappa$ in (a) is the unitary exponent of Theorem~\ref{thm:main-group}(c): a service with purity $B$ is read at relator defect about
$(B/n)^{1/2}$. The lower bound in (b) needs no assumption on purity, because at the rank rate purity is paid for with memory
(Proposition~\ref{prop:purity-budget}): there $P+C<3Q+6$, and the argument of (a) with $B$ of order $Q$ gives $Q^2\gtrsim n$ up to logarithms. The
services in (a) and (b) are the two devices of Section~\ref{sec:devices}. We prove Theorem~\ref{thm:main-channel} in
Section~\ref{sec:main-channel-proof}, after the purity budget and the upper bound, and read it as a statement about finite tracial
approximations in Corollary~\ref{cor:profile}.
